% Context from chapters-tex/wavelets.tex; for mathematical reference, not compilation.
\section{Multiresolution Detail Spaces}

The detail spaces are the orthogonal complements associated with the inclusions $V_j \subset V_{j-1}$. Since these subspaces are closed,
\eq{
	\foralls j, \quad W_j \text{ is such that } V_{j-1} = V_j \oplus^\bot W_j. 
}
The following diagram shows these nested approximation spaces and their orthogonal complements:
\begin{center}
	\image{wavelets}{.8}{embeded-spaces-1d}
\end{center}
Suppose that $W_0$ admits an orthonormal basis of translates $\{\psi(\cdot-n)\}_n$. Scaling then gives
\eq{
	\{ \psi_{j,n} \}_n \text{ is a Hilbertian orthonormal basis of } W_j
	\qwhereq
	\psi_{j,n} \eqdef \frac{1}{2^{j/2}} \psi\pa{ \frac{\cdot - 2^j n}{2^j} }.
}

Orthogonal decomposition across all scales yields
\eq{
	L^2(\RR) = \bigoplus_{j=-\infty}^{j=+\infty} W_j
	=V_{j_0}\oplus^\bot\bigoplus_{j\leq j_0}W_j.
}
For every $f \in L^2(\RR)$, the corresponding expansion converges in $L^2(\RR)$:
\eq{
	f = 
	\lim_{(j_-,j_+) \rightarrow (-\infty,+\infty)} \sum_{j=j_-}^{j_+} P_{W_j} f
	= 
	\lim_{j_-\to-\infty}\left(P_{V_{j_0}}f+\sum_{j=j_-}^{j_0}P_{W_j}f\right).
}
This decomposition shows that 
\eq{
	\enscond{\psi_{j,n}}{(j,n) \in \ZZ^2}
}
is an orthonormal basis of $L^2(\RR)$, which is called a wavelet basis.
 
Stopping at a coarsest scale gives the orthonormal basis
\eq{
	\enscond{\psi_{j,n}}{j \leq j_0, n \in \ZZ} \cup 
	\enscond{\phi_{j_0,n}}{n \in \ZZ}.
}

The forward wavelet transform computes all inner products of a function $f$ with the elements of this basis.

\myfigure{
\tabdeux{
\image{wavelets}{.48}{multires-1d-haar-0}&
\image{wavelets}{.48}{multires-1d-haar-1}\\
Signal $f$ & $j=-8$ \\
\image{wavelets}{.48}{multires-1d-haar-2}&
\image{wavelets}{.48}{multires-1d-haar-3}\\
$j=-7$ & $j=-6$
}
}{ 
	1-D Haar multiresolution projection $P_{V_j}f$ of a function $f$.  
}{fig-haar-1d}



  
\paragraph{Haar wavelets.}

For the Haar multiresolution~\eqref{eq-haar-multires}, one has
\eql{\label{eq-haar-details}
	W_j = \enscond{ f }{  \foralls n \in \ZZ, \; f \text{ constant on  } [2^{j-1}n,2^{j-1}(n+1)) \qandq
	 \int_{n2^j}^{(n+1)2^j} f = 0 }.
}
One choice of mother wavelet is
\eq{
	\psi(t)=\choice{
		1 \qforq 0 \leq t < 1/2, \\
		-1 \qforq 1/2 \leq t < 1, \\
		0 \quad \text{otherwise,}
	}
}
as shown in Figure~\ref{fig-haard-scaling-wav}.


\myfigure{
\tabdeux{
\image{wavelets}{.48}{multires-1d-haar-1}& \\
$P_{V_j} f$ & \\
\image{wavelets}{.48}{multires-1d-haar-2}&
\image{wavelets}{.48}{multires-1d-haar-detail-1}\\
$P_{V_{j+1}} f$ & $P_{W_{j+1}} f$ \\
\image{wavelets}{.48}{multires-1d-haar-3}&
\image{wavelets}{.48}{multires-1d-haar-detail-2} \\
$P_{V_{j+2}} f$ & $P_{W_{j+2}} f$
}
}{ 
	Successive Haar approximations (left) and the extracted details (right). Increasing the relative level from $j$ to $j+1$ and $j+2$ coarsens the approximation. 
}{fig-haar-appr-details}

Figure~\ref{fig-haar-appr-details} illustrates how subtracting consecutive approximation-space projections isolates the details at each scale.


  
\paragraph{Shannon and splines.}

Figure~\ref{fig-spline-vs-shannon} compares the sharp frequency partitions of Shannon wavelets with the smooth transitions of spline wavelets. Shannon wavelets can be viewed as a limiting case as the spline degree increases.


\begin{figure}[tbp]
\centering
\BookDrawing{tikz/wavelets-shannon-spline-bands}
\caption{Spline and Shannon wavelet segmentation of the frequency axis.}
\label{fig-spline-vs-shannon}
\end{figure}



                                                   
\section{On Bounded Domains}

On the periodic domain $\TT=\RR/\ZZ$, we use period 1 rather than the $2\pi$-period convention of the Fourier chapter. To obtain an orthonormal wavelet basis of $L^2(\TT)$, periodize the wavelets and restrict their translation points $2^j n$ to $[0,1]$, choosing $0 \leq n < 2^{-j}$. As in~\eqref{eq-periodizin